%%%
\section{1. Introducing GL Controller}

\begin{frame}{1.0: GL Chamber: Introduction}{} 
GL Chamber will eventually replace all P300-based (and/or SCP-based) chambers. The Platinous Series based on GL controller has been released (as we speak). Work on Global chamber is underway. 
\begin{figure} [hbt!] 
\centering
\includegraphics[scale=0.28]{figures/glc-chamber-001a2.png}
\caption{ESPEC Platinous Series (EPL, EPU, EPX, EPX models, J-series)}
\end{figure}
\end{frame}

\begin{frame}{1.1: What about Benchtop Chamber?}{} 
Benchtop chamber still uses Watlow F4T. Reason: Benchtop E-box is too small to fit GL controller hardware (I/O board, etc). For now, Benchtop chambers will be offered along with GL chambers. 
\begin{figure} [hbt!] 
\centering
\includegraphics[scale=0.34]{figures/benchtop-f4t.jpg}
\caption{E-box of Benchtop with ESPEC Web Controller}
\end{figure}
\end{frame}

\begin{frame}{1.2: Good News \& Bad News?}{}

\begin{minipage}[c]{0.50\textwidth}
{\bf{Good News?}} 
 We still offer our legacy ESPEC Web Controller (EWC) system for Benchtop chambers and U-Web kit for existing ESPEC (or Non-ESPEC with F4T) chambers that Customers still own. We can still apply what we know about EWC to support customers.    
\end{minipage} 
\hfill
\begin{minipage}[c]{0.450\textwidth} 
        \includegraphics[width=\textwidth]{figures/benchtop-chambers-002.png} 
        \captionof{figure}{Benchtops with F4T}
\end{minipage}

\vspace{2em}
\begin{minipage}[c]{1.00\textwidth}
{\bf{Bad News?}} 
 We now have to learn about the GL controller and GL chambers to help potential customers. {\it{Good news or bad news, it's your view. This short training will hopefully change all of that to ``good news.''}}  The ``good news'' part: GL controller in many ways is similar to our EWC. {\it{Details are explained in the next several slides.}}  
\end{minipage}
\hfill
%\begin{minipage}[c]{0.450\textwidth} 
%        \includegraphics[width=\textwidth]{figures/up2board-io-002.PNG} 
%        \captionof{figure}{Ethernet ports on Up2board}
%\end{minipage}
\end{frame}

\section{2. What is GL Controller?}

\begin{frame}{2.0: GL Controller: What is GL Chamber?}{}
\begin{minipage}[c]{0.50\textwidth}
GL controller is part hardware and part software that makes up the ``GL chamber." 

\vspace{0.1in} 
The GL hardware part makes up the Input/Output board (see figure).

\vspace{0.1in}
The software part is written in CODESYS; its details do not concern us! The software part works with our ``legacy EWC software'' that runs on the EWC hardware, called Up2Board (figure below). 
\begin{figure} [hbt!] 
      \centering
      \includegraphics[scale=.20]{figures/up2board-001a.PNG}  
      %\caption{EWC V3 log-in screen on a Web browser}
\end{figure}
\end{minipage} 
\hfill
\begin{minipage}[c]{0.450\textwidth} 
        \includegraphics[width=\textwidth]{figures/platinous-jseries-001.png} 
        \captionof{figure}{GL controller software and hardware in the E-box of Platinous J Series} 
\end{minipage}

\end{frame} 

\begin{frame}{2.1: GL System vs. Legacy EWC}{}
\begin{minipage}[c]{0.520\textwidth}
Figure on the right depicts how a Benchtop chamber and EWC work together to provide control and operation of the chamber. Chamber interface and communication between F4T and EWC are via Ethernet or serial (RS232) connection. Depicted in the figure, two Ethernet cables run from EWC, one to F4T, the other to external Ethernet port to join the main network for remote access of the EWC user interface. 

\vspace{1em}
In contrast, GL system and EWC are integrated, to be described in detail below. 

\end{minipage} 
\hfill
\begin{minipage}[c]{0.4250\textwidth} 
        \includegraphics[width=\textwidth]{figures/benchtop-f4t.jpg}
        \captionof{figure}{EWC and Benchtop F4T} 
\end{minipage}
\end{frame} 

\begin{frame}{2.2: Details of EWC and P300 Interface (Review)}{}
\begin{minipage}[c]{0.530\textwidth}
Figure on the right depicts communication interface between EWC and P300. We already know, EWC connects to a P300 controller via a serial interface to provide control and operation of the chamber. 

\vspace{1em}
The same communication protocol is also used on SCP-220 (or ES102) or Watlow F4. 

\vspace{1em}
{\bf{Summary:}} Legacy EWC functions as a separate device to connect to a controller (such as P300, SCP220, F4/F4T, etc) to provide control and operation of the chamber. 
\end{minipage} 
\hfill
\begin{minipage}[c]{0.420\textwidth} 
        \includegraphics[width=\textwidth]{figures/ewc-n-p300.PNG}
        \captionof{figure}{GL Chamber} 
\end{minipage}
\end{frame} 

\begin{frame}{2.3: Details of GL \& EWC Interface}{}
\begin{minipage}[c]{0.530\textwidth}

Unlike the legacy EWC system which offers chamber interface and configuration as an external connection (i.e. because it is implemented to support various controllers), the GL controller has no external chamber interface. 

\vspace{1em} GL controller {\it{is}} the CODESYS that resides on the EWC platform. EWC and CODESYS are the GL controller. The GL controller uses Up2Board (EWC hardware) to link to its Input/Output board. 
\end{minipage} 
\hfill
\begin{minipage}[c]{0.420\textwidth} 
        \includegraphics[width=\textwidth]{figures/glc-chamber-003.png}
        \captionof{figure}{GL Chamber} 
\end{minipage}
\end{frame} 

 
\begin{frame}{2.4: Detailed Interface of GL Controller}{}
\begin{minipage}[c]{0.530\textwidth}

Interface \& I/O ports of EWC hardware used by GL controller.

    \begin{figure} [hbt!] 
      \centering
      \includegraphics[scale=0.250]{figures/gl-up2board-001a.png}
      \caption{Interfaces used by GL Ctrl} 
    \end{figure}
    
    \begin{itemize} 
        \item [(1)] Supplies power to HMI 
        \item [(2)] Connects to HMI operation 
        \item [(3)] links CODESYS license; fourth USB for external storage
        \item [(4)] Connects to HDMI display 
        \item [(5)] Connects to main network  
    \end{itemize} 
\end{minipage} 
\hfill
\begin{minipage}[c]{0.420\textwidth} 
        \includegraphics[width=\textwidth]{figures/glc-chamber-004.png}
        \captionof{figure}{GL Chamber} 
\end{minipage}
\end{frame} 

\begin{frame}{2.5: Chamber User Interface (HMI)}{}
Complete control and operation of the chamber are accomplished via the dedicated (touchscreen) HMI on the chamber.  

    \begin{figure} [hbt!] 
      \centering
      \includegraphics[scale=0.250]{figures/gl-hmi-003a.png}
      \caption{GL Touchscreen Display, dedicated HMI} 
    \end{figure}
\end{frame}   

\begin{frame}{2.6: Details of External Ports of GL Chamber}{}

Depicted in the figure below, external ports of the GL controller has three ports: (1) Product Temp. Protector, (2) Chamber light switch, (3) External Storage (USB Thumb drive for downloading data, uploading update packages, etc). 

    \begin{figure} [hbt!] 
      \centering
      \includegraphics[scale=0.40]{figures/gl-external-ports-001a.png}
      \caption{External ports on GL chamber} 
    \end{figure}
\end{frame}   

% SECTION 3
\section{3. GL Software Update and Distribution} 
\begin{frame}{3.0: GL Software Update}{}
Similar to our current (legacy) EWC, GL system uses chamber S/N as its registration number to receive firmware update through our cloud service provider, MenderIO.
    \begin{figure} [hbt!] 
      \centering
%      \includegraphics[scale=.30]{figures/menderio-cloud-001.PNG}  
      \includegraphics[scale=0.300]{figures/mender-cloud-service-002.PNG}
      \caption{GL System subscribes to MenderIO cloud service}
    \end{figure}
\end{frame} 

\begin{frame}{3.1: Firmware Update: Online or Offline}{} 
Customer has complete control over firmware update method: Offline or Online through a cloud service. 
\begin{enumerate}
\item Using a stepladder technique, update package is pushed onto the next available root partition (see diagram). 

    \begin{figure} [hbt!] 
      \centering
%      \includegraphics[scale=.30]{figures/menderio-cloud-001.PNG}  
      \includegraphics[scale=.60]{figures/stepladder-001.PNG}
      \caption{Stepladder technique of GL update, (same as EWC)}
    \end{figure}

\item If the update process is successful, the system will boot and run on this new firmware.   
\item If the update process failed, the system rolls back to the previous known operation state. 
\end{enumerate} 
\end{frame} 

\begin{frame}{3.2: Firmware Update Subscription}{} 
Firmware update is provided through a yearly update subscription. Customers purchase the subscription 
to receive the updates. 
\begin{itemize}
   \item {\bf{Online Update}}: \\
       For GL chamber that connects to the Internet, MenderIO cloud updating service handles the update process. 
       Automatic update will take place when GL chamber is in Standby mode.   
   \item {\bf{Offline Update}}:   \\
       For GL chamber that does not connect to the Internet, the owner is provided with a link to access their account
       to download an update package to perform the update procedure on their system. Refer to User's
       Manual for detail on Firmware Update procedure. 
\end{itemize} 
\end{frame} 

% SECTION 4
% Troubleshooting 
\section{4. User Interface: Display Options, Features \& Operation}

\begin{frame}{4.0: GL Operation Panel}{}

The operating panel is the 10-inch touchscreen HMI that provides the User Interface (UI) for control and operation of the chamber. 
\begin{figure} [hbt!] 
      \centering
      \includegraphics[scale=.30]{figures/gl-ui-001.PNG}  
      \caption{User Interface of GL Controller}
\end{figure}
The UI consists of three major components: (1) Menu bar, (2) Status bar, and (3) Standard Display area. The touchscreen is operated by gently pressing the screen elements provided by these components.
\end{frame} 

\begin{frame}{4.1: GL Main Menu \& Features}{} 

\begin{enumerate}
   \item {\bf{Home}}:Home page of GL controller. If any menu button is pressed/clicked, the login dialog appears for the user to log in.
   \item {\bf{Operation Mode}}: If a user is not logged in, operation buttons for Constant or Program are locked and grayed out.
   \item {\bf{Monitor}}: This menu displays the current status of the GL controller.
   \item {\bf{Set Mode}}: If a user is not logged in, contents of Constant Setup or Program Setup are viewable in read-only mode; they cannot be modified.
   \item {\bf{Setup}}: If a user is not logged in, only few submenus are accessible; the rest are locked out.
\end{enumerate} 

Hands-on will provide a better understanding of these menus, many of which are similar to our current ESPEC Web Controller menu and user interface. 

\end{frame}

\begin{frame}{4.2: Three Ways to Access GL Menu}{}

GL Controller system offers three ways to access its system menu. 

\begin{figure} [hbt!] 
      \centering
      \includegraphics[scale=.3250]{figures/three-access-way-001.PNG}  
      \caption{User Interface of GL Controller}
\end{figure}

\begin{itemize} 
    \item[(1)] Menu bar 
    \item[(2)] Menu Button (via the Home page)
    \item[(3)] Hamburger Button (in Status bar) 
\end{itemize} 
\end{frame} 


% SECTION 5

\section{5. GL Chamber \& Network Configuration} 
\begin{frame}{5.0: DHCP \& Static Nework Configuration}{}
By default, the GL controller operates as a standalone system. Complete control and operation of the instrumentation are achieved via the dedicated (touchscreen) HMI on the chamber. 

\vspace{1em}
When the GL controller is connected to a network, you can use a Web browser to access the GL controller user interface (UI) to control and operate the chamber. The GL controller utilizes a standard Web browser to provide (i.e., host) its UI for control and operation of the instrumentation. 

\vspace{1em}
Thanks to the ability of a Web browser to operate on any computer (or device) on the network, GL controller operations can be performed remotely by authorized users on the network. {\it{Same old stuff we know from EWC.}} 
\end{frame}


\begin{frame}{5.1: Network Configuration} {} 
\begin{itemize} 
   \item DHCP: IP address handled by DHCP server
       \begin{figure} [hbt!] 
         \centering
         \includegraphics[scale=.14]{figures/1672855610-dhcp-setup-figure.PNG}  
         \caption{DHCP network diagram}
       \end{figure}

   \item Static Network: IP address configured manually
      \begin{figure} [hbt!] 
        \centering
        \includegraphics[scale=.18]{figures/2789777267-typhoon-static-setup-v2.PNG}  
        \caption{Static network diagram}
      \end{figure}
\end{itemize} 
\end{frame}

\begin{frame}{5.2: GL Chamber: Standalone}{}

GL Chamber can operate by itself without connecting to a network or an Internet. It can still operate as a standalone system even when it is part of a network.  
\begin{figure} [hbt!] 
      \centering
      \includegraphics[scale=.30]{figures/glc-chamber-001a2.png}
      \caption{GL a standalone system}
\end{figure}
\end{frame} 

\begin{frame}{5.3: GL Chamber: Network Communication}{}

As part of a network, GL chamber can be accessed remotely using a Web browser via IP address or hostname of the GL chamber. However, remote access must be enabled for this to happen. It must be enabled on the HMI via the Setup submenu ``Set Protection''. This option is linked to remote communication in the next slide. 

\begin{figure} [hbt!] 
      \centering
      \includegraphics[scale=0.40]{figures/gl-network-001.png}  
      \caption{GL chamber as part of a network device}
\end{figure}
\end{frame} 


\section{6. GL Controller \& ChamberConnectLibrary}

\begin{frame}{6.0: GL Remote Communication}{} 
Communication protocol of the GL controller system is a general-purpose communication tool that allows communication between the GL chamber and remote device (PC). Communication protocol is currently available for TCP and RS232 interface, with TCP set by default. This option allows customers to control the GL chamber remotely (or locally) using GL native text commands. 

\vspace{1em} 

\begin{itemize} 
  \item Locally: This feature is available under the Set Communication submenu: {\keys{Setup}}$\longrightarrow${\keys{Configuration}}$\longrightarrow${\keys{Set Communication}}
  \item Remotely: Remote communication can be enabled/disabled. This option is available under the Set Protection submenu: {\keys{Setup}}$\longrightarrow${\keys{Set Protection}} or on HMI Set Protection button.   
\end{itemize}
\end{frame}

\begin{frame}{6.1: GL \& ChamberConnectLibrary}{} 
Chamber Connect Library using Python 3 will be available for customers who prefer to customize program to control their GL chamber remotely. The Library will be integrated with the current Chamber Connect Library for P300, SCP220, ES102 and Watlow F4 and F4T.  

\end{frame}

\section{7. Troubleshooting GL Chamber} 


\begin{frame}{7.0: Troubleshooting Issues on GL}{}

  Integration of hardware \& software of the GL controller has a number of advantages. Issue customer may have is virtually not about chamber misconfiguration or chamber connection interface any more. How many times have we helped customers with their chamber interface connection and misconfiguration that they themselves did.  

  \begin{enumerate} 
\item Server Error /Chamber error: Stopped Wroking
\item Unable to access GL chamber remotely
\item Unable to issue set commands directly to GL chamber
\item Network connection issues 
  \end{enumerate} 
\end{frame}  

\section{8. Universal Web Controller Kit for GL Chamber?} 

\begin{frame}{8.0: U-Web Kit Support: COMM Interface}{} 

The Universal Web Server (Controller) kit continues with its standard configuration options to provide connection, operation and control on existing ESPEC or Non-ESPEC (with F4T) chambers. We stil ldo offer this option. Default communication with F4T is via TCP/IP; with other chamber/controllers, it will use the serial COMM via the Serial-to-USB interface, as shown below.  

    \begin{figure} [hbt!] 
      \centering
      \includegraphics[scale=.220]{figures/webkit0.PNG}  \\
      \caption{U Web kit hardware in box and bare}
    \end{figure}
\end{frame}


\begin{frame}{8.1: U-Web Kit for Legacy Chambers}{} 
EWC can be sold as a Universal Web Kit to connect and control the existing chambers (mentioned above). 

    \begin{figure} [hbt!] 
      \centering
      \includegraphics[scale=.260]{figures/u-web-kit-up2board-001.PNG}  
      \caption{U-Web Kit with USB-to-Serial interface cable}
    \end{figure}
{\bf{Advantages:}} 
\begin{itemize}
\item U-Web kit is an excellent option for adding EWC to the existing chamber for control and operation. 
\item Setup is generally made simple with its quick guide and preconfiguration. 
\end{itemize} 
\end{frame}

\begin{frame}{8.2: U-Web Kit for GL Chambers}{} 
  {\center {\large { 
There is no U-Web kit for GL chambers. GL {\it{is}} the Web Controller. 
}
 }}


\end{frame}

\section{9. Questions/Answers \& Hands-On Demo} 

\begin{frame}{9.0: Q/A \& GL Hands-on}{}
\begin{itemize}
   \item Questions \& Answers
    \item Quick demo of GL chamber operation and features

    \begin{figure} [hbt!] 
      \centering
      \includegraphics[scale=.150]{figures/gl-chamber-op1.PNG}  
      \caption{Program/Preview menu of GL controller}
    \end{figure}

    \begin{figure} [hbt!] 
      \centering
      \includegraphics[scale=.150]{figures/gl-chamber-op2.PNG}  
      \caption{Program execution and monitoring on GL chamber}
    \end{figure}

\end{itemize}
\end{frame}